% GENERATED FILE. Edit canonical lessons, then run npm run generate:books.
% canonical-source-hash: fnv1a64:df6b80f776555dff
% canonical-lessons: KA-C546-kaitota, KA-C546-nerehore, KA-C546-sunna, KA-C546-tarasi, KA-C546-udu, KA-R546-review-540-546, KA-R546-review-final-546

\chapter{Home Comforts}
\label{ch:ka-home-comforts}
\hlchaptermodality{546}

\begin{chapteropening}
\textbf{By the end of this chapter:} \emph{I can say a kitchen garden, a neighbourhood, lime, a terrace and to puff.}

Complete the last lesson of chapter 546: I can say a kitchen garden, a neighbourhood, lime, a terrace and to puff.
\end{chapteropening}

\section[\texorpdfstring{\kn{ಕೈತೋಟ} (kaitōṭa)}{ಕೈತೋಟ (kaitōṭa)}]{\kn{ಕೈತೋಟ} (kaitōṭa) — a kitchen garden}

\label{lesson:KA-C546-kaitota}



\begin{quote}
Before the new one: say the Kannada for to wash clothes, then the Kannada for starch.
\end{quote}

\subsection*{You'll want to know: \kn{ಕೈತೋಟ}}
\textbf{\kn{ಕೈತೋಟ}} — \emph{kaitōṭa} — \textquotedblleft{}a kitchen garden\textquotedblright{}.

\begin{quote}
\kn{ನಮ್ಮ} \kn{ಕೈತೋಟದಲ್ಲಿ} \kn{ಟೊಮೇಟೊ} \kn{ಬೆಳೆಯುತ್ತದೆ}.
\end{quote}

\begin{quote}
\emph{namma kaitōṭadalli ṭomēṭo beḷeyuttade.}
\end{quote}

\begin{quote}
Tomatoes grow in our kitchen garden.
\end{quote}

\kn{ಕೈ}, hand + \kn{ತೋಟ}, garden: a garden small enough to tend by hand.

\begin{cousinweb}
\textbf{\kn{ಕೈತೋಟ}} (\emph{kaitōṭa}) is a Kannada compound: native \emph{kai}, \textquotedblleft{}hand\textquotedblright{} (DEDR 2023), and \emph{tōṭa}, \textquotedblleft{}a garden\textquotedblright{}.
\end{cousinweb}

\begin{grammarlens}[title={What you've built}]
One more word for describing what you lived and felt.
\end{grammarlens}

\subsection*{Your turn}
\begin{itemize}
\raggedright
  \item \emph{Say it:} \emph{kaitōṭa}
  \item \emph{Say it:} \emph{kaitōṭa}, once more
  \item \emph{Say it:} \emph{namma kaitōṭadalli ṭomēṭo beḷeyuttade.} — Tomatoes grow in our kitchen garden.
  \item \emph{Recall:} say the Kannada for to wash clothes, then the Kannada for starch, then say \emph{kaitōṭa} again
\end{itemize}

\begin{tcolorbox}[breakable,colback=teal!4,colframe=teal!35!black,title={Before you move on}]
What does \kn{ಕೈತೋಟ} mean? (\textquotedblleft{}A kitchen garden\textquotedblright{}.) Say it once more.
\end{tcolorbox}
\section[\texorpdfstring{\kn{ನೆರೆಹೊರೆ} (nerehore)}{ನೆರೆಹೊರೆ (nerehore)}]{\kn{ನೆರೆಹೊರೆ} (nerehore) — a neighbourhood}

\label{lesson:KA-C546-nerehore}



\begin{quote}
Before the new one: say the Kannada for starch, then the Kannada for a kitchen garden.
\end{quote}

\subsection*{You'll want to know: \kn{ನೆರೆಹೊರೆ}}
\textbf{\kn{ನೆರೆಹೊರೆ}} — \emph{nerehore} — \textquotedblleft{}a neighbourhood\textquotedblright{}.

\begin{quote}
\kn{ನಮ್ಮ} \kn{ನೆರೆಹೊರೆಯಲ್ಲಿ} \kn{ಎಲ್ಲರೂ} \kn{ಸ್ನೇಹದಿಂದ} \kn{ಇದ್ದಾರೆ}.
\end{quote}

\begin{quote}
\emph{namma nerehoreyalli ellarū snēhadinda iddāre.}
\end{quote}

\begin{quote}
Everyone in our neighbourhood is friendly.
\end{quote}

\kn{ನೆರೆಯವರು} are the neighbours themselves.

\begin{cousinweb}
\textbf{\kn{ನೆರೆಹೊರೆ}} (\emph{nerehore}) is a native Kannada pair-word built on \emph{nere}, \textquotedblleft{}adjoining, proximity, neighbourhood\textquotedblright{} (DEDR 3770); the second half echoes it.
\end{cousinweb}

\begin{grammarlens}[title={What you've built}]
One more word for describing what you lived and felt.
\end{grammarlens}

\subsection*{Your turn}
\begin{itemize}
\raggedright
  \item \emph{Say it:} \emph{nerehore}
  \item \emph{Say it:} \emph{nerehore}, once more
  \item \emph{Say it:} \emph{namma nerehoreyalli ellarū snēhadinda iddāre.} — Everyone in our neighbourhood is friendly.
  \item \emph{Recall:} say the Kannada for starch, then the Kannada for a kitchen garden, then say \emph{nerehore} again
\end{itemize}

\begin{tcolorbox}[breakable,colback=teal!4,colframe=teal!35!black,title={Before you move on}]
What does \kn{ನೆರೆಹೊರೆ} mean? (\textquotedblleft{}A neighbourhood\textquotedblright{}.) Say it once more.
\end{tcolorbox}
\section[\texorpdfstring{\kn{ಸುಣ್ಣ} (suṇṇa)}{ಸುಣ್ಣ (suṇṇa)}]{\kn{ಸುಣ್ಣ} (suṇṇa) — lime; whitewash}

\label{lesson:KA-C546-sunna}



\begin{quote}
Before the new one: say the Kannada for a kitchen garden, then the Kannada for a neighbourhood.
\end{quote}

\subsection*{You'll want to know: \kn{ಸುಣ್ಣ}}
\textbf{\kn{ಸುಣ್ಣ}} — \emph{suṇṇa} — \textquotedblleft{}lime; whitewash\textquotedblright{}.

\begin{quote}
\kn{ಹಬ್ಬಕ್ಕೆ} \kn{ಮುಂಚೆ} \kn{ನಾವು} \kn{ಗೋಡೆಗಳಿಗೆ} \kn{ಸುಣ್ಣ} \kn{ಬಳಿದೆವು}.
\end{quote}

\begin{quote}
\emph{habbakke muñce nāvu gōḍegaḷige suṇṇa baḷidevu.}
\end{quote}

\begin{quote}
Before the festival we whitewashed the walls.
\end{quote}

\kn{ಸುಣ್ಣ} \kn{ಬಳಿ}, to whitewash. \kn{ಸೀಮೆಸುಣ್ಣ}, chalk, holds the same word.

\begin{cousinweb}
\textbf{\kn{ಸುಣ್ಣ}} (\emph{suṇṇa}) goes back to Sanskrit \emph{cūrṇa}, \textquotedblleft{}something ground, a powder\textquotedblright{} (CDIAL 4889), through a Prakrit form.
\end{cousinweb}

\begin{grammarlens}[title={What you've built}]
One more word for describing what you lived and felt.
\end{grammarlens}

\subsection*{Your turn}
\begin{itemize}
\raggedright
  \item \emph{Say it:} \emph{suṇṇa}
  \item \emph{Say it:} \emph{suṇṇa}, once more
  \item \emph{Say it:} \emph{habbakke muñce nāvu gōḍegaḷige suṇṇa baḷidevu.} — Before the festival we whitewashed the walls.
  \item \emph{Recall:} say the Kannada for a kitchen garden, then the Kannada for a neighbourhood, then say \emph{suṇṇa} again
\end{itemize}

\begin{tcolorbox}[breakable,colback=teal!4,colframe=teal!35!black,title={Before you move on}]
What does \kn{ಸುಣ್ಣ} mean? (\textquotedblleft{}Lime; whitewash\textquotedblright{}.) Say it once more.
\end{tcolorbox}
\section[\texorpdfstring{\kn{ತಾರಸಿ} (tārasi)}{ತಾರಸಿ (tārasi)}]{\kn{ತಾರಸಿ} (tārasi) — a terrace; a flat concrete roof}

\label{lesson:KA-C546-tarasi}



\begin{quote}
Before the new one: say the Kannada for a neighbourhood, then the Kannada for lime.
\end{quote}

\subsection*{You'll want to know: \kn{ತಾರಸಿ}}
\textbf{\kn{ತಾರಸಿ}} — \emph{tārasi} — \textquotedblleft{}a terrace; a flat concrete roof\textquotedblright{}.

\begin{quote}
\kn{ಬೇಸಿಗೆಯಲ್ಲಿ} \kn{ನಾವು} \kn{ತಾರಸಿಯ} \kn{ಮೇಲೆ} \kn{ಮಲಗುತ್ತೇವೆ}.
\end{quote}

\begin{quote}
\emph{bēsigeyalli nāvu tārasiya mēle malaguttēve.}
\end{quote}

\begin{quote}
In summer we sleep on the terrace.
\end{quote}

\kn{ತಾರಸಿ} \kn{ಮನೆ}, a house with a flat concrete roof.

\begin{cousinweb}
\textbf{\kn{ತಾರಸಿ}} (\emph{tārasi}) is the Kannada word for a flat concrete roof; its origin is not settled.
\end{cousinweb}

\begin{grammarlens}[title={What you've built}]
One more word for describing what you lived and felt.
\end{grammarlens}

\subsection*{Your turn}
\begin{itemize}
\raggedright
  \item \emph{Say it:} \emph{tārasi}
  \item \emph{Say it:} \emph{tārasi}, once more
  \item \emph{Say it:} \emph{bēsigeyalli nāvu tārasiya mēle malaguttēve.} — In summer we sleep on the terrace.
  \item \emph{Recall:} say the Kannada for a neighbourhood, then the Kannada for lime, then say \emph{tārasi} again
\end{itemize}

\begin{tcolorbox}[breakable,colback=teal!4,colframe=teal!35!black,title={Before you move on}]
What does \kn{ತಾರಸಿ} mean? (\textquotedblleft{}A terrace; a flat concrete roof\textquotedblright{}.) Say it once more.
\end{tcolorbox}
\section[\texorpdfstring{\kn{ಊದು} (ūdu)}{ಊದು (ūdu)}]{\kn{ಊದು} (ūdu) — to puff; to blow into}

\label{lesson:KA-C546-udu}



\begin{quote}
Before the new one: say the Kannada for lime, then the Kannada for a terrace.
\end{quote}

\subsection*{You'll want to know: \kn{ಊದು}}
\textbf{\kn{ಊದು}} — \emph{ūdu} — \textquotedblleft{}to puff; to blow into\textquotedblright{}.

\begin{quote}
\kn{ಅಜ್ಜಿ} \kn{ಒಲೆಯ} \kn{ಬೆಂಕಿಯನ್ನು} \kn{ಕೊಳವೆಯಿಂದ} \kn{ಊದಿದರು}.
\end{quote}

\begin{quote}
\emph{ajji oleya beṅkiyannu koḷaveyinda ūdidaru.}
\end{quote}

\begin{quote}
Grandmother blew on the stove fire through a pipe.
\end{quote}

Past (I …): \kn{ಊದಿದೆ}, \emph{ūdide}.

With the mouth; the wind's blowing is \kn{ಬೀಸು}.

\begin{cousinweb}
\textbf{\kn{ಊದು}} (\emph{ūdu}) is native Dravidian: DEDR 741 lists Kannada \emph{ūdu}, \textquotedblleft{}to blow, to become inflated, swell\textquotedblright{}.
\end{cousinweb}

\begin{grammarlens}[title={What you've built}]
One more word for describing what you lived and felt.
\end{grammarlens}

\subsection*{Your turn}
\begin{itemize}
\raggedright
  \item \emph{Say it:} \emph{ūdu}
  \item \emph{Say it:} \emph{ūdu}, once more
  \item \emph{Say it:} \emph{ajji oleya beṅkiyannu koḷaveyinda ūdidaru.} — Grandmother blew on the stove fire through a pipe.
  \item \emph{Recall:} say the Kannada for lime, then the Kannada for a terrace, then say \emph{ūdu} again
\end{itemize}

\begin{tcolorbox}[breakable,colback=teal!4,colframe=teal!35!black,title={Before you move on}]
What does \kn{ಊದು} mean? (\textquotedblleft{}To puff; to blow into\textquotedblright{}.) Say it once more.
\end{tcolorbox}
\section[(dialogue)]{Looking back: sport, fitness and the home}

\label{lesson:KA-R546-review-540-546}



\begin{quote}
Say the Kannada for kabaddi and wrestling.
\end{quote}

\subsection*{Your turn}
Say these aloud:

\begin{itemize}
\raggedright
  \item kabaddi, wrestling, a goal in sport, to kick, to punch
  \item a workout, endurance, mountaineering, archery, to cheer on
  \item an upper storey, a threshold, a hall, a basement, a cabinet
  \item a veranda, a lantern, a cushion, furniture, to lean back
  \item a switch, a plug, a gas cylinder, a tube, a wire
  \item a rug, a quilt, to shake out, to wash clothes, starch
  \item a kitchen garden, a neighbourhood, lime, a terrace, to puff
\end{itemize}

\begin{tcolorbox}[breakable,colback=teal!4,colframe=teal!35!black,title={Before you move on}]
Kabaddi? (\textbf{\kn{ಕಬಡ್ಡಿ}}, \emph{kabaḍḍi}.) Wrestling? (\textbf{\kn{ಕುಸ್ತಿ}}, \emph{kusti}.) A goal in sport? (\textbf{\kn{ಗೋಲು}}, \emph{gōlu}.) To kick? (\textbf{\kn{ಒದೆ}}, \emph{ode}.) To punch? (\textbf{\kn{ಗುದ್ದು}}, \emph{guddu}.) A workout? (\textbf{\kn{ಕಸರತ್ತು}}, \emph{kasarattu}.) Endurance? (\textbf{\kn{ಸಹಿಷ್ಣುತೆ}}, \emph{sahiṣṇute}.) Mountaineering? (\textbf{\kn{ಪರ್ವತಾರೋಹಣ}}, \emph{parvatārōhaṇa}.) Archery? (\textbf{\kn{ಬಿಲ್ಲುಗಾರಿಕೆ}}, \emph{billugārike}.) To cheer on? (\textbf{\kn{ಹುರಿದುಂಬಿಸು}}, \emph{huridumbisu}.) An upper storey? (\textbf{\kn{ಮಾಳಿಗೆ}}, \emph{māḷige}.) A threshold? (\textbf{\kn{ಹೊಸ್ತಿಲು}}, \emph{hostilu}.) A hall? (\textbf{\kn{ಹಜಾರ}}, \emph{hajāra}.) A basement? (\textbf{\kn{ನೆಲಮಾಳಿಗೆ}}, \emph{nelamāḷige}.) A cabinet? (\textbf{\kn{ಕಪಾಟು}}, \emph{kapāṭu}.) A veranda? (\textbf{\kn{ವರಾಂಡ}}, \emph{varāṇḍa}.) A lantern? (\textbf{\kn{ಲಾಂದ್ರ}}, \emph{lāndra}.) A cushion? (\textbf{\kn{ಮೆತ್ತೆ}}, \emph{mette}.) Furniture? (\textbf{\kn{ಪೀಠೋಪಕರಣ}}, \emph{pīṭhōpakaraṇa}.) To lean back? (\textbf{\kn{ಒರಗು}}, \emph{oragu}.) A switch? (\textbf{\kn{ಸ್ವಿಚ್}}, \emph{svic}.) A plug? (\textbf{\kn{ಪ್ಲಗ್}}, \emph{plag}.) A gas cylinder? (\textbf{\kn{ಸಿಲಿಂಡರ್}}, \emph{siliṇḍar}.) A tube? (\textbf{\kn{ಕೊಳವೆ}}, \emph{koḷave}.) A wire? (\textbf{\kn{ತಂತಿ}}, \emph{tanti}.) A rug? (\textbf{\kn{ಜಮಖಾನ}}, \emph{jamakhāna}.) A quilt? (\textbf{\kn{ರಗ್ಗು}}, \emph{raggu}.) To shake out? (\textbf{\kn{ಝಾಡಿಸು}}, \emph{jhāḍisu}.) To wash clothes? (\textbf{\kn{ಒಗೆ}}, \emph{oge}.) Starch? (\textbf{\kn{ಗಂಜಿ}}, \emph{gañji}.) A kitchen garden? (\textbf{\kn{ಕೈತೋಟ}}, \emph{kaitōṭa}.) A neighbourhood? (\textbf{\kn{ನೆರೆಹೊರೆ}}, \emph{nerehore}.) Lime? (\textbf{\kn{ಸುಣ್ಣ}}, \emph{suṇṇa}.) A terrace? (\textbf{\kn{ತಾರಸಿ}}, \emph{tārasi}.) To puff? (\textbf{\kn{ಊದು}}, \emph{ūdu}.)
\end{tcolorbox}
\section[(dialogue)]{Home Comforts — once more}

\label{lesson:KA-R546-review-final-546}



\begin{quote}
Say the Kannada for a kitchen garden and a neighbourhood.
\end{quote}

\subsection*{Your turn}
\begin{itemize}
\raggedright
  \item \emph{Say it:} a kitchen garden, a neighbourhood, lime, a terrace, to puff
\end{itemize}

\begin{tcolorbox}[breakable,colback=teal!4,colframe=teal!35!black,title={Before you move on}]
A kitchen garden? (\textbf{\kn{ಕೈತೋಟ}}, \emph{kaitōṭa}.) A neighbourhood? (\textbf{\kn{ನೆರೆಹೊರೆ}}, \emph{nerehore}.) Lime? (\textbf{\kn{ಸುಣ್ಣ}}, \emph{suṇṇa}.) A terrace? (\textbf{\kn{ತಾರಸಿ}}, \emph{tārasi}.) To puff? (\textbf{\kn{ಊದು}}, \emph{ūdu}.)
\end{tcolorbox}
